\documentclass[11pt]{report}
\usepackage[utf8]{inputenc}
\usepackage{amssymb,amsthm,mathtools}
\usepackage{hyperref,graphics}
\usepackage{bigints}
\usepackage{todonotes}
\usepackage{natbib}

\usetikzlibrary{decorations.markings}
\graphicspath{{images/}}

\newcommand{\eps}{\varepsilon}
\newcommand{\hsp}{\hspace{0.3cm}}

\delimitershortfall-1sp
\newcommand{\field}[1]{\mathbb{#1}}

\newcommand{\abs}[1]{\left|#1\right|}
\newcommand{\norm}[1]{\left\lVert#1\right\rVert}

\newcommand{\cbr}[1]{\left(#1\right)}
\newcommand{\bbr}[1]{\left[#1\right]}
\newcommand{\cbbr}[1]{\left(#1\right]}
\newcommand{\bcbr}[1]{\left[#1\right)}
\newcommand{\fbr}[1]{\left\{#1\right\}}

\newcommand{\mysetminusD}{\hbox{\tikz{\draw[line width=0.6pt,line cap=round] (3pt,0) -- (0,6pt);}}}
\newcommand{\mysetminusT}{\mysetminusD}
\newcommand{\mysetminusS}{\hbox{\tikz{\draw[line width=0.45pt,line cap=round] (2pt,0) -- (0,4pt);}}}
\newcommand{\mysetminusSS}{\hbox{\tikz{\draw[line width=0.4pt,line cap=round] (1.5pt,0) -- (0,3pt);}}}
\newcommand{\mysetminus}{\mathbin{\mathchoice{\mysetminusD}{\mysetminusT}{\mysetminusS}{\mysetminusSS}}}

\DeclareMathOperator{\dist}{\text{dist}}

\newtheorem{definition}{Definition}[chapter]
\newtheorem{theorem}{Theorem}[section]
\newtheorem{corollary}{Corollary}[theorem]
\newtheorem{lemma}[theorem]{Lemma}
\newtheorem*{remark}{Remark}

\begin{document}

\title{
	{Riemann Zeta Function - The Functional Equation}\\
	{\large University of Vienna}\\
	{\vspace{3ex}\includegraphics[width=15ex]{University.png}}
}
\author{Florian Jesacher}
\date{April 1, 2016}
\maketitle

\begin{abstract}
We first introduce the Riemann zeta function in the right half-plane and use Weierstrass's uniform convergence theorem to address the holomorphy of the function. By deriving its Euler product representation we deduce the absence of roots on the function's domain. Once we have defined the gamma function and examined some of its identities including Euler's reflection formula, we create a link to the Riemann zeta function. We manage to meromorphically continue said function and extend it to the whole complex plane by utilising the analytic character of contour integrals. This leads us directly to a meromorphic integral expression with only one simple pole finally paving the way to the functional equation. One fundamental result of this new expression is that the Riemann zeta function has no nontrivial zeros outside the critical strip.
\end{abstract}

\tableofcontents

\chapter{Introduction}
The Riemann zeta function is one of the most important functions in analytic number theory, particularly since it encodes a lot of information about prime numbers. Thus it is well worth taking a closer look at its properties and behaviour in the complex plane. Special attention should be paid to the zeros of the function as they might give some indication of the mysterious distribution of primes. \\
This work is mainly based on \cite{Apostol1976} and uses some definitions and theorems from \cite{Ahlfors1966} as well. \\
Henceforth let $s \in \field{C}$ be a complex number where $\sigma \coloneqq \Re(s)$ and $t \coloneqq \Im(s)$. The space of all holomorphic functions on a domain $S \subseteq \field{C}$ shall be denoted by $\mathcal{H}(S) \coloneqq \fbr{f \in \field{C} ^S \colon f \textit{ holomorphic}}$.


\section{Elementary Properties of $\zeta$}
First of all we want to draw some conclusions about the convergence of the Riemann zeta function and its holomorphy on its domain.


\begin{definition}
	We refer to
\begin{equation*}
\begin{aligned}	
	&\zeta \colon \fbr{s \in \field{C} \colon \sigma > 1} \to \field{C}, \\
	&\zeta(s) = \sum_{n = 1}^{\infty} n^{-s}
\end{aligned}
\end{equation*}
	as the Riemann zeta function.
\end{definition}


\subsection{Absolute Convergence and Holomorphy}


\begin{theorem}
	$\zeta(s)$ converges absolutely for $\sigma > 1$.
\end{theorem}
\begin{proof}
	Let $x \in \field{R}_{>0}$. We quickly note that
\begin{equation*}
	\abs{x^s}
	= \abs{e^{s \log(x)}}
	= \abs{e^{(\sigma + i t) \log(x)}}
	= \abs{x^{\sigma} e^{i t \log(x)}}
	= x^{\sigma}
\end{equation*}
	as
\begin{equation*}
	\forall \theta \in \field{R} \colon \abs{e^{i \theta}} = 1.
\end{equation*}
	Hence we obtain
\begin{equation*}
	\sum_{n = 1}^{\infty} \abs{n^{-s}} = \sum_{n = 1}^{\infty} n^{-\sigma}.
\end{equation*}
	We know that the right sum converges for $\sigma > 1$. Consequently $\zeta(s)$ converges absolutely.
\end{proof}


\begin{lemma}[Weierstrass's uniform convergence theorem]
	Let $S \subseteq \field{C}$ be open, $f \in \field{C}^S$ and $(f_n)_{n \in \field{N}} \in \mathcal{H}(S)^{\field{N}}$ with $\lim\limits_{n \to \infty} f_n = f$. If
\begin{equation*}
	\forall D \subset S \textit{ compact}: f_n\big\vert _D \to f\big\vert _D \textit{ uniformly},
\end{equation*}
	then
\begin{equation}
	f \in \mathcal{H}(S)
\end{equation}
	and
\begin{equation}
	f_n^{(k)}\big\vert _D \to f^{(k)}\big\vert _D \textit{ uniformly for } k \in \field{N}.
\end{equation}
\end{lemma}
\begin{proof}
	Let $\triangle \subset S$ denote a closed solid triangle. As $\triangle$ is compact and due to the uniform convergence of $f_n\big\vert _\triangle \to f\big\vert _\triangle$ and Goursat’s integral lemma we know that
\begin{equation*}
	\forall \triangle \subset S \colon \int_{\partial \triangle} f(z) dz = \int_{\partial \triangle} \lim\limits_{n \to \infty} f_n(z) dz	= \lim\limits_{n \to \infty}\int_{\partial \triangle} f_n(z) dz = 0.
\end{equation*}
	 Morera's theorem finally gives us
\begin{equation*}
	f \in \mathcal{H}(S).
\end{equation*}
	To show the second claim, let $r \in \field{R}_{> 0}$ be the radius of an arbitrary closed disk $D_r \subset S$ and let the positively orientated boundary of $D_r$ be denoted by $\partial D_r$. Since $f_n^{(k)} \in \mathcal{H}(S)$ we can apply Cauchy's integral formula
\begin{equation*}
\begin{aligned}
	f^{(k)}(z)
	&= \frac{k!}{2 \pi i} \int _{\partial D_r} \frac{f(\xi)}{\cbr{\xi - z}^{k + 1}} d\xi \\
	&= \frac{k!}{2 \pi i} \int _{\partial D_r} \lim\limits_{n \to \infty} \frac{f_n(\xi)}{\cbr{\xi - z}^{k + 1}} d\xi \\
	&= \lim\limits_{n \to \infty} \frac{k!}{2 \pi i} \int _{\partial D_r} \frac{f_n(\xi)}{\cbr{\xi - z}^{k + 1}} d\xi
	= \lim\limits_{n \to \infty} f_n^{(k)}(z)
\end{aligned}
\end{equation*}
	for $z \in D_r^\circ$. For $l < r$ we obtain
\begin{equation*}
\begin{aligned}
	\norm{f_n^{(k)} - f^{(k)}}_{D_l}
	&\leq \sup\limits_{z \in D_l} \fbr{\frac{k!}{2 \pi} \int _{\partial D_r}\abs{\frac{f_n(\xi) - f(\xi)}{\cbr{\xi - z}^{k + 1}}} \abs{d\xi}} \\
	&\leq \sup\limits_{z \in D_l} \fbr{\frac{k! r}{\dist(z, \partial D_r)^{k + 1}} \norm{f_n - f}_{\partial D_r}} \\
	&\leq \frac{k! r}{(r - l)^{k + 1}} \norm{f_n - f}_{\partial D_r}.
\end{aligned}
\end{equation*}
	Hence $f_n^{(k)}\big\vert _{D_l} \to f^{(k)}\big\vert _{D_l}$ uniformly. Since we can always find a finite cover of open disks $D^i$ for any compact set $D \subset S$ such that $\bigcup \overline{D^i} \subset S$, we have provided conclusive evidence for our claim.
\end{proof}


\begin{theorem}
	$\zeta$ is holomorphic for $\sigma > 1$.
\end{theorem}
\begin{proof}
	Let $R = \bbr{a, b} \times i\bbr{c, d} \subset \field{C}$ with $1 < \delta \leq a$. By the definition of uniform convergence we write
\begin{equation*}
\begin{aligned}
	&\forall \eps > 0 \; \exists N \in \field{N} \; \forall m \geq N \; \forall s \in R \colon \\
	&\begin{aligned}
	\abs{\sum _{n = 1} ^\infty n^{-s} - \sum _{n = 1} ^{m - 1} n^{-s}}
	= \abs{\sum _{n = m} ^\infty n^{-s}}
	&\leq \sum _{n = m} ^\infty \abs{n^{-s}}
	= \sum _{n = m} ^\infty n^{-\sigma} \\
	&\leq \sum _{n = m} ^\infty n^{-\delta} < \eps.
	\end{aligned}
\end{aligned}
\end{equation*}
	Therefore $\zeta_m \big\vert _R \to \zeta \big\vert _R$ uniformly and according to Weierstrass $\zeta$ is holomorphic for $\sigma > 1$.
\end{proof}


\section{The Euler Product}
In this section we introduce infinite products and talk about their convergence. After deriving $\zeta$'s Euler product representation we can make some assertions about the occurrence of zeros for $\sigma > 1$.


\subsection{Infinite Convergent Products}


\begin{definition}
	Let $(a_n)_{n \in \field{N}} \in \field{C}^\field{N}$. The infinite product
\begin{equation*}
	\prod_{n \in \field{N}} \cbr{1 + a_n}
\end{equation*}
	is said to be convergent if
\begin{equation*}
	\lim\limits_{N \to \infty} \prod _{n = 1} ^N \cbr{1 + a_n} \in \field{C} \mysetminus \fbr{0}.
\end{equation*}
	We speak of absolute convergence if the same holds true for $\abs{a_n}$.
\end{definition}


\begin{remark}
	We particularly don't allow convergent products to be zero here, since any infinite product being zero would converge if only one factor was zero, hence the convergence wouldn't depend on the whole sequence of factors. Furthermore this convention would be too radical for certain applications.
\end{remark}


\begin{lemma}
	Let $(a_n)_{n \in \field{N}} \in \field{C}^\field{N}$. The following statements are equivalent
\begin{equation*}
\begin{aligned}
	&(i) &&\prod_{n \in \field{N}} \cbr{1 + a_n} \textit{ converges} \\
	&(ii) &&\sum_{n \in \field{N}} \log(1 + a_n) \textit{ converges}.
\end{aligned}
\end{equation*}
\end{lemma}
\begin{proof}
	Let the infinite product be denoted by $P$, its partial products by $P_n$ and respectively the series by $S$ and its partial sums by $S_n$. If $P_n$ converges then
\begin{equation*}
	\lim\limits_{n \to \infty} \frac{P_n}{P_{n - 1}} = \lim\limits_{n \to \infty} \cbr{1 + a_n} = 1.
\end{equation*}
	Hence $(a_n)_{n \in \field{N}}$ has to be a null sequence. To prove $(i) \implies (ii)$ we assume that $P_n \to P$ and we choose the principal value of $\log(P)$, namely
\begin{equation*}
	\log(P) = \log(\abs{P}) + i \arg(P).
\end{equation*}
	Now we determine $\arg(P_n)$ by the condition
\begin{equation*}	
	\arg(P) - \pi < \arg(P_n) \leq \arg(P) + \pi
\end{equation*}
	and set $\log(P_n) = \log(\abs{P_n}) + i \arg(P_n)$. Since $\log$ is a multivalued function we get
\begin{equation*}
	S_n = \log(P_n) + h_n 2 \pi i
\end{equation*}
	for $h_n \in \field{Z}$. For two consecutive terms we can write
\begin{equation*}
	(h_{n + 1} - h_n) 2 \pi i = \log(1 + a_{n + 1}) + \log(P_n) - \log(P_{n + 1}).
\end{equation*}
	If we choose $n$ large enough we get
\begin{equation*}
\begin{aligned}
	\abs{\arg(1 + a_{n + 1})} &< \frac{2 \pi}{3} \\
	\abs{\arg(P) - \arg(P_n)} &< \frac{2 \pi}{3} \\
	\abs{\arg(P) - \arg(P_{n + 1})} &< \frac{2 \pi}{3}.
\end{aligned}
\end{equation*}
	These inequalities imply
\begin{equation*}
	\abs{h_{n + 1} - h_n} < 1,
\end{equation*}
	hence for a sufficiently large $n$ we have $h_n = h_{n + 1} = h$ and have
\begin{equation*}
	\lim\limits_{n \to \infty} S_n = S = \log(P) + h 2 \pi i.
\end{equation*}
	To show $(ii) \implies (i)$ we assume that $S_n \to S$ and by writing $P_n = e^{S_n}$ for the principal branch of $\log$ we conclude that
\begin{equation*}
	\lim\limits_{n \to \infty} P_n = P = e^{S} \in \field{C} \mysetminus {0}.
\end{equation*}
\end{proof}


\begin{lemma}
	Let $(a_n)_{n \in \field{N}} \in \field{C}^\field{N}$. The following statements are equivalent
\begin{equation*}
\begin{aligned}
	&(i) &&\prod_{n \in \field{N}} \cbr{1 + \abs{a_n}} \textit{ converges} \\
	&(ii) &&\sum_{n \in \field{N}} \abs{a_n} \textit{ converges}.
\end{aligned}
\end{equation*}
\end{lemma}
\begin{proof}
	We can write
\begin{equation*}
	0 < 1 + \sum_{n \in \field{N}} \abs{a_n} \leq \prod_{n \in \field{N}} \cbr{1 + \abs{a_n}} = e^{\sum_{n \in \field{N}} \log(1 + \abs{a_n})} \leq e^{\sum_{n \in \field{N}} \abs{a_n}}
\end{equation*}
	as
\begin{equation*}
	\forall x \in \field{R}_{\geq 0} \colon \log(1 + x) \leq x.
\end{equation*}
\end{proof}


\subsection{Euler Product Representation of $\zeta$}


\begin{definition}
	An arithmetic function $f \in \field{C}^\field{N}$ satisfying
\begin{equation*}
	\forall n,m \in \field{N} \colon \gcd\cbr{n, m} = 1 \implies f(n m) = f(n) f(m)
\end{equation*}
	is called multiplicative and we particularly have $f(1) = 1$. If even
\begin{equation*}
	\forall n,m \in \field{N} \colon f(n m) = f(n) f(m)
\end{equation*}
	holds true we say $f$ is completely multiplicative.
\end{definition}


\begin{theorem}
	Let $\mathcal{P} = \fbr{p \in \field{N} \colon \; p \textit{ prime number}}$ and let $f \in \field{C}^\field{N}$ be an arithmetic function. If $\sum_{n \in \field{N}} f(n)$ converges absolutely and if $f$ is multiplicative, the series can be expressed by an absolutely convergent infinite product
\begin{equation*}
	\sum_{n \in \field{N}} f(n) = \prod_{p \in \mathcal{P}} \cbr{1 + \sum_{i = 1}^\infty f(p^i)}.
\end{equation*}
	If $f$ even is completely multiplicative and fulfils $\abs{f(p)} < 1 \; \forall p \in \mathcal{P}$ we have
\begin{equation*}
	\sum_{n \in \field{N}} f(n) = \prod_{p \in \mathcal{P}} \frac{1}{1 - f(p)}.
\end{equation*}
\end{theorem}
\begin{proof}
	Let $\mathcal{P}_{\leq x} \coloneqq \fbr{p \in \mathcal{P} \colon \; p \leq x}$. Take the product
\begin{equation*}
	Q(x) = \prod_{p \in \mathcal{P}_{\leq x}} \cbr{1 + \sum_{i = 1}^\infty f(p^i)}.
\end{equation*}
	Since the product is finite and our series converge absolutely we can expand the product and rearrange its terms without altering the sum. For an arbitrary multiplication of two factors involving primes $p_k$ and $p_l$ we get
\begin{equation*}
\begin{aligned}		
	\cbr{1 + \sum_{i = 1}^\infty f(p_k^i)} \cbr{1 + \sum_{j = 1}^\infty f(p_l^j)} \\
	= 1 + \sum_{i = 1}^\infty f(p_k^i) + \sum_{j = 1}^\infty \cbr{f(p_l^j) + \sum_{i = 1}^\infty f(p_k^i) f(p_l^j)}
\end{aligned}
\end{equation*}
	and because $f$ is multiplicative, we observe that after calculating the whole product we end up with summands of the form
\begin{equation*}
	f(p_{i_1}^{\alpha_1})f(p_{i_2}^{\alpha_2}) \dots f(p_{i_r}^{\alpha_r}) = f(p_{i_1}^{\alpha_1} p_{i_2}^{\alpha_2} \dots p_{i_r}^{\alpha_r}),
\end{equation*}
	where $p_{i_1}^{\alpha_1} p_{i_2}^{\alpha_2} \dots p_{i_r}^{\alpha_r}$ occurs only once within the whole sum as the above representation is unique due to the unique factorisation theorem. This allows us to write
\begin{equation*}
	Q(x) = \sum_{n \in A} f(n),
\end{equation*}
	where $A$ consists of those $n$ having all their prime factors $\leq x$ and $1 \in A$, hence
\begin{equation*}
	\sum_{n \in \field{N}} f(n)	 - Q(x) = \sum_{n \in B} f(n),
\end{equation*}
	where B is the set of $n$ having at least one prime factor $> x$. Thus
\begin{equation*}
	\abs{\sum_{n \in \field{N}} f(n) - Q(x)} \leq \sum_{n \in B} \abs{f(n)} \leq \sum_{n > x} \abs{f(n)}
\end{equation*}
	converges to $0$ as $x \to \infty$ implying
\begin{equation*}	
	\lim\limits_{x \to \infty} Q(x) = \sum_{n \in \field{N}} f(n).
\end{equation*}
	To show the absolute convergence of $Q(\infty)$ we recall that the infinite product $\prod_{n \in \field{N}} \cbr{1 + a_n}$ converges absolutely if and only if $\sum_{n \in \field{N}} \abs{a_n} < \infty$. We write
\begin{equation*}
	\forall x \in \field{N} \colon \; \sum_{p \in \mathcal{P}_{\leq x}} \abs{\sum_{i = 1}^\infty f(p^i)} \leq \sum_{p \in \mathcal{P}_{\leq x}} \sum_{i = 1}^\infty \abs{f(p^i)} \leq \sum_{n = 2}^\infty \abs{f(n)} < \infty
\end{equation*}
	which implies
\begin{equation*}
	\sum_{p \in \mathcal{P}} \abs{\sum_{i = 1}^\infty f(p^i)} < \infty.
\end{equation*}
	Therefore $Q(\infty)$ is absolutely convergent. If $f$ is completely multiplicative and fulfils $\abs{f(p)} < 1 \; \forall p \in \mathcal{P}$ the infinite product simplifies to
\begin{equation*}
	\sum_{n \in \field{N}} f(n) = \prod_{p \in \mathcal{P}} \cbr{1 + \sum_{i = 1}^\infty f(p^i)} = \prod_{p \in \mathcal{P}} \cbr{\sum_{i = 0}^\infty f(p)^i} = \prod_{p \in \mathcal{P}} \frac{1}{1 - f(p)}.
\end{equation*}
\end{proof}


\begin{corollary}
	We have $\zeta$'s Euler product representation
\begin{equation*}
	\zeta(s) = \prod_{p \in \mathcal{P}} \frac{1}{1 - p^{-s}} \hsp \textit{ for } \sigma > 1.
\end{equation*}
\end{corollary}
\begin{proof}
	Since $f(n) = n^{-s}$ is a completely multiplicative arithmetic function, $\abs{f(n)} = n^{-\sigma} < 1$ for $\sigma > 1$ and the series converges absolutely the above statement follows directly from the previous theorem.
\end{proof}


\begin{corollary}
	$\zeta(s)$ has no zeros for $\sigma > 1$.
\end{corollary}
\begin{proof}	
	According to the last theorem we know that the infinite product representation of $\zeta(s)$ for $\sigma > 1$ is absolutely convergent, which implies that $\zeta(s)$ can't vanish for $\sigma > 1$. 
\end{proof}

\chapter{Meromorphic Continuation}
Analytic continuation provides a way to extend the domain of a given analytic complex function. We can expand this function into a power series which is only valid within its radius of convergence. Under certain circumstances these expansions have a larger-than-expected radius of convergence and its power series can be used to define the function outside of its domain. If said function has at most countable many poles we speak of meromorphic continuation. We won't explicitly use power series  for our meromorphic continuation of the Riemann zeta function but a ruse where contour integrals are involved.


\section{The Gamma Function}
The gamma function is closely related to the Riemann zeta function and plays an important role in $\zeta$'s meromorphic continuation beyond the $\sigma = 1$ line. Thus we want to acquire some of its properties including Euler's reflection formula, which helps us to make assertions about zero-free regions of the Riemann zeta function.


\begin{definition}
	The function
\begin{equation*}
\begin{aligned}
	&\Gamma \colon \fbr{s \in \field{C} \colon \sigma > 0} \to \field{C}, \\
	&\Gamma(s) = \int _0 ^\infty x^{s - 1} e^{-x} dx
\end{aligned}
\end{equation*}
	is called gamma function.
\end{definition}


\subsection{Convergence and Holomorphy}


\begin{theorem}
	$\Gamma(s)$ is well defined for $\sigma > 0$
\end{theorem}
\begin{proof}
	We can write
\begin{equation*}
\begin{aligned}
	\abs{\Gamma(s)}
	&= \abs{\int _0 ^\infty x^{s - 1} e^{-x} dx} \leq \int _0 ^\infty \abs{x^{s - 1} e^{-x}} \\
	&= \int _0 ^1 x^{\sigma - 1} e^{-x} dx + \int _1 ^k x^{\sigma - 1} e^{-x} dx + \int _k ^\infty x^{\sigma - 1} e^{-x} dx
\end{aligned}
\end{equation*}
	for $k \in \cbr{1, \infty}$. For the first integral we find
\begin{equation*}
	\int _0 ^1 x^{\sigma - 1} e^{-x} dx \leq \int _0 ^1 x^{\sigma - 1} dx = \frac{1}{\sigma}
\end{equation*}
	since $e^{-x} \leq 1$ for $x > 0$. As the second integral clearly exists for every $k$, we proceed to the third integral. Note that
\begin{equation*}
	\forall r \in \field{R} \colon \lim\limits_{x \to \infty} x^r e^{-\frac{x}{2}} = 0
\end{equation*}
 	hence
\begin{equation*}
	\exists k \in \cbr{0, \infty} \; \forall x > 0 \geq k \colon \; x^r e^{-\frac{x}{2}} \leq 1.
\end{equation*}
	Thus we have
\begin{equation*}
	\int _k ^\infty x^{\sigma - 1} e^{-x} dx = \int _k ^\infty \cbr{x^{\sigma - 1} e^{-\frac{x}{2}}} e^{-\frac{x}{2}} dx \leq \int _k ^\infty e^{-\frac{x}{2}} dx = 2 e^{-\frac{k}{2}}.
\end{equation*}
\end{proof}


\begin{theorem}
	$\Gamma(s)$ is holomorphic for $\sigma > 0$.
\end{theorem}
\begin{proof}
	Let $R = \bbr{a, b} \times i\bbr{c, d} \subset \field{C}$ with $0 < \delta \leq a$. From the previous theorem we deduce
\begin{equation*}
\begin{aligned}
	&\forall \eps > 0 \; \exists k \in \field{N} \; \forall m \geq k \; \forall s \in R \colon \\
	&\abs{\int _0 ^\infty x^{s - 1} e^{-x} dx - \int _0 ^m x^{s - 1} e^{-x} dx} \leq \abs{\int _m ^\infty x^{s - 1} e^{-x} dx} \leq 2 e^{-\frac{m}{2}} < \eps,
\end{aligned}
\end{equation*}
	hence $\Gamma_m \big\vert _R \to \Gamma \big\vert _R$ uniformly and according to Weierstrass $\Gamma$ is holomorphic for $\sigma > 0$.
\end{proof}


\begin{theorem}
	We have the identity
\begin{equation*}
	\Gamma(s + 1) = s \Gamma(s).
\end{equation*}
\end{theorem}
\begin{proof}
	By applying integration by parts we get
\begin{equation*}
	\Gamma(s + 1) = \int _0 ^\infty x^{s} e^{-x} dx =\bbr{-x^s e^{-x}} _0 ^\infty + s \int _0 ^\infty x^{s - 1} e^{-x} dx = s\Gamma(s).
\end{equation*}
\end{proof}


\subsection{Meromorphic Continuation}


\begin{theorem}
	$\Gamma(s)$ is meromorphically continuable for all $s \in \field{C}$ with simple poles at $s = -n\in \field{Z} _{\leq 0}$ and we have
\begin{equation*}
	Res\cbr{\Gamma(s), -n} = \frac{\cbr{-1}^n}{n!}.
\end{equation*}
\end{theorem}
\begin{proof}
	For fixed $s \in \field{C} \mysetminus \field{Z}_{\leq 0}$ let $n \in \field{N}_0$ be such that $\sigma + n > 0$. Then we have
\begin{equation*}
	\Gamma(s) = \frac{1}{s} \Gamma(s + 1) = \dots = \frac{1}{s \cbr{s + 1} \dots \cbr{s + n - 1}} \Gamma(s + n) = \frac{\Gamma(s + n)}{P_n(s)}.
\end{equation*}
	Thus $\Gamma$'s holomorphy in the negative plane is reducible to the holomorphy of the function in the positive plane which we've already shown. Hence $\Gamma(s)$ is well defined and its poles are just the simple roots of $\lim\limits_{n \to \infty} P_n(s)$, namely $\ker\cbr{P_\infty(s)} = \field{Z}_{\leq 0}$. \\
	Now we want to calculate the residues of $\Gamma(s)$ and find
\begin{equation*}
	Res\cbr{\Gamma(s), -n} = \lim\limits_{s \to -n} \cbr{s + n} \frac{\Gamma(s + n + 1)}{s \cbr{s + 1} \dots \cbr{s + n}} = \frac{\cbr{-1}^n}{n!}
\end{equation*}
\end{proof}


\begin{lemma}
	In particular we have
\begin{equation*}
	\Gamma(n + 1) = n! \textit{ for } n \in \field{N}_0.
\end{equation*}
\end{lemma}
\begin{proof}
	We easily check that
\begin{equation*}
	\Gamma(n + 1) = n \Gamma(n) = \dots = n! \Gamma(1) = n!.
\end{equation*}
\end{proof}


\subsection{Euler's Reflection Formula}


\begin{lemma}
	It holds that
\begin{equation*}
	\forall \sigma \in \field{R} \colon \lim\limits_{\abs{t} \to \infty} \Gamma(\sigma + i t) = 0,
\end{equation*}
	so $\Gamma(\sigma + i t)$ is bounded for a fixed $\sigma \in \field{R}$ and $\abs{t}$ large enough.
\end{lemma}
\begin{proof}
	Let $\sigma >0$ be fixed. For the substitution $y = \frac{1}{2 \pi} \log(x)$ we get
\begin{equation*}
\begin{aligned}	
	\Gamma(\sigma + i t)
	&= \int_0^\infty x^{\sigma + i t - 1} e^{-x} dx \\
	&= \int_{-\infty}^\infty e^{2 \pi i y t} e^{2 \pi y (\sigma - 1)} e^{-e^{2 \pi y}} dy = \int_{-\infty}^\infty e^{2 \pi i y t} \hat{g}(y) dy.
\end{aligned}
\end{equation*}
	Per definition of the Fourier transform, the above formula says that the function
\begin{equation*}
\begin{aligned}
	&g\colon \field{R} \to \field{C}, \\
	&g(t) = \Gamma(\sigma + i t)
\end{aligned}
\end{equation*}
	is the inverse Fourier transform of $\hat{g}$. Since the (inverse) Fourier transform describes an isomorphism on the Schwartz space $\mathcal{S}$, we deduce that $g$ has to decrease more rapidly for an increasing $\abs{t}$ than any polynomial and therefore $g \in \mathcal{S}(\field{R})$. In other words it in fact follows that
\begin{equation*}
	\lim\limits_{\abs{t} \to \infty} g(t) = \lim\limits_{\abs{t} \to \infty} \Gamma(\sigma + i t) = 0.
\end{equation*}
	By meromorphic continuation of $\Gamma$ this holds for all $\sigma \in \field{R}$.
\end{proof}


\begin{theorem}[Euler's reflection formula]
	We have
\begin{equation*}
	\Gamma(s)\Gamma(1 - s) = \frac{\pi}{\sin(\pi s)} \hsp \textit{  for } s \in \field{C} \mysetminus \field{Z}.
\end{equation*}
	Hence $\Gamma(s)$ has no zeros and its reciprocal function $\frac{1}{\Gamma(s)}$ is well defined with zeros at the simple poles of $\Gamma(s)$.
\end{theorem}
\begin{proof}
	Let us consider
\begin{equation*}
	f(s) = \sin(\pi s) \Gamma(s) \Gamma(1 - s).
\end{equation*}
	We want to show that $f(s)$ is constant. It is easy to check that $f(s)$ is entire as
\begin{equation*}
	\lim\limits_{s \to 0} \frac{\sin(s)}{s} = 1,
\end{equation*}
	hence $\sin(s)$ eradicates the simple poles of $\Gamma(s)$. Furthermore $f(s)$ satisfies
\begin{equation*}
\begin{aligned}
	f(s + 1) 
	&= \sin(\pi \cbr{s + 1}) \Gamma(s + 1) \Gamma(-s) \\ 
	&= -\sin(\pi s) s \Gamma(s) \cbr{-\frac{1}{s}} \Gamma(1 - s) = f(s).
\end{aligned}
\end{equation*}
	So $f_t(\sigma + it)$ is $1$-periodic for each fixed $t$ and $f_t \in C^\infty(\field{R} \mysetminus \field{Z})$. Therefore $f(s)$ has a rapidly decreasing converging Fourier series
\begin{equation*}
	f(s) = \sum _{n \in \field{Z}} c_n(t) e^{i 2 \pi n \sigma}.
\end{equation*}
	Let us consider the $m$-th term of the Fourier series for fixed $s$. By definition we have
\begin{equation*}
	f_m(s) = \int _0 ^1 f(s + x) e^{- i 2 \pi m x} dx = c_m(t) e^{i 2 \pi m \sigma}.
\end{equation*}
	Since $f_m(s)$ is an analytic function, the Cauchy-Riemann equations imply that
\begin{equation*}
	\frac{d}{dt}\cbr{c_m(t) \cos(2 \pi m \sigma)} = - \frac{d}{d\sigma}\cbr{2 \pi m c_m(t) \sin(2 \pi m \sigma)},
\end{equation*}
	or equivalently
\begin{equation*}
	\frac{d}{dt} c_m(t) = -2 \pi m c_m(t),
\end{equation*}
	which represents a linear homogenous differential equation with the general solution
\begin{equation*}
	c_m(t) = c_m e^{-2 \pi m t}.
\end{equation*}
	Finally we can write
\begin{equation*}
	f(s) = \sum _{m \in \field{Z}} c_m e^{i 2 \pi m s}.
\end{equation*}
	Let us now suppose $m > 0$ and $t \to -\infty$ or if $m < 0$ let $t \to \infty$. Then we have
\begin{equation*}
	f_m(\sigma + i t) = c_m e^{-i 2 \pi m \sigma} e^{2 \pi m \abs{t}} = \widetilde{c}_m e^{2 \pi m \abs{t}}.
\end{equation*}
	On the other hand
\begin{equation*}
	\abs{f_m(\sigma + i t)} = \abs{\int _0 ^1 f(\sigma + x + i t) e^{- i 2 \pi m x} dx} \leq \int _0 ^1 \abs{f(\sigma + x + i t)} dx.
\end{equation*}
	According to previous lemma  $\abs{\Gamma(\sigma + i t)}$ is bounded for $\abs{t}$ large enough and we get
\begin{equation*}
\begin{aligned}
	\abs{\widetilde{c}_m} e^{2 \pi m \abs{t}}
	&\leq \int_0^1 \abs{\sin(x + \sigma + i t)} \abs{\Gamma(x + \sigma + i t) \Gamma(x + 1 - \sigma - i t)} dx \\
	&\leq \int_0^1 c e^{\pi \abs{t}} dx
	\leq d e^{\pi \abs{t}}
\end{aligned}
\end{equation*}
	where $c,d \in \field{R}_{> 0}$. Thus by comparing the growth rates we conclude that $\widetilde{c}_m = 0$ for $m \neq 0$ leaving $f(s)$ constant. The value is easily determined by
\begin{equation*}
\begin{aligned}
	\lim\limits_{s \to 0} f(s)
	&= \lim\limits_{s \to 0} \sin(s) \Gamma(s) \Gamma(1 - s) \\
	&= \lim\limits_{s \to 0} \frac{\sin(\pi s)}{s} \Gamma(s + 1) \Gamma(1 - s)
	= \pi.
\end{aligned}	
\end{equation*}

\end{proof}


\section{Contour Integral Representation}
After deriving an integral representation of the product of the Gamma and the Riemann zeta function, we manage to meromorphically continue $\zeta$ beyond the $\sigma = 1$ line. We use the fact that path integrals over continuous functions along a part-wise continuous path $\lambda$ are analytic on the whole $\field{C} \mysetminus \abs{\lambda}$.


\begin{theorem}
	We have the identity
\begin{equation}
	\zeta(s)\Gamma(s) = \int _0 ^\infty \frac{x^{s - 1} e^{-x}}{1 - e^{-x}} dx \textit{ for } \sigma > 1.
\end{equation}
\end{theorem}
\begin{proof}
	First we assume that $s \in \field{R}_{> 1}$. Let $x = n t$ where $n \in \field{N}_{> 0}$. By substitution we get
\begin{equation*}
	\Gamma(s) = \int _0 ^{\infty} x^{s - 1} e^{-x} (1 - e^{-x})^{-1} dx = n^s \int _0 ^\infty e^{-n t} t^{s - 1} dt,
\end{equation*}
equivalently
\begin{equation*}
	n^{-s} \Gamma(s) = \int _0 ^\infty e^{-n t} t^{s - 1} dt.
\end{equation*}
	After summing up both sides over $n$ ranging from $1$ to $\infty$ we obtain
\begin{equation*}
	\zeta(s) \Gamma(s) = \sum _{n=1} ^\infty n^{-s} \Gamma(s) = \sum _{n=0} ^\infty \int _0 ^\infty e^{-(n + 1) t} t^{s - 1} dt.
\end{equation*}
	Since the integral can be treated as a Lebesgue integral and since the integrand is nonnegative, Levi’s convergence theorem tells us that the series
\begin{equation*}
	\sum _{n=0} ^\infty e^{-(n + 1) t} t^{s - 1} dt
\end{equation*}
	converges almost everywhere to a Lebesgue-integrable sum function on $\bcbr{O, \infty}$ and that we can interchange sum and integral such that
\begin{equation*}
	\zeta(s) \Gamma(s) = \int _0 ^\infty \sum _{n = 0} ^\infty e^{-(n + 1) t} t^{s - 1} dt.
\end{equation*}
	If $t > 0$ then $0 < e^{-t} < 1$. Thus
\begin{equation*}
	\sum _{n = 0} ^\infty e^{n t} = \frac{1}{1 - e^{-t}}
\end{equation*}
	is a geometric series. Bringing it together we have
\begin{equation*}
	\sum _{n=0} ^\infty e^{-(n + 1) t} t^{s - 1} = \frac{e^{-t} t^{s - 1}}{1 - e^{-t}}
\end{equation*}
	everywhere on $\bcbr{0, \infty}$, except for $s = 0$, so we get
\begin{equation*}
	\zeta(s) \Gamma(s) = \int _0 ^\infty \frac{e^{-t} t^{s - 1}}{1 - e^{-t}} dt
\end{equation*}
	for real $s > 1$. To extend it to all complex $s$ with $\sigma > 1$ we note that both members are analytic for $\sigma > 1$. To show that the right member is analytic we assume $1 + \delta \leq \sigma \leq c$, where $c > 1$ and $\delta > 0$ and write
\begin{equation*}
	\int _0 ^\infty \abs{\frac{e^{-t} t^{s - 1}}{1 - e^{-t}}} dt \leq \int _0 ^\infty \frac{e^{-t} t^{\sigma - 1}}{1 - e^{-t}} dt = \cbr{\int _0 ^1 + \int _1 ^\infty} \frac{e^{-t} t^{\sigma - 1}}{1 - e^{-t}} dt.
\end{equation*}
	If $0 \leq t \leq 1$ we have $t^{\sigma - 1} \leq t^\delta$ and if $t \geq 1$ we have $t^{\sigma - 1} \leq t^{c - 1}$. Also, since $e^t - 1 \geq t$ for $t > 0$ we have
\begin{equation*}
	\int _0 ^1 \frac{e^{-t} t^{\sigma - 1}}{1 - e^{-t}} dt \leq \int _0 ^1 \frac{t^\delta}{e^t - 1} \leq \int _0 ^1 t^{\delta - 1} = \frac{1}{\delta}
\end{equation*}
	and
\begin{equation*}
	\int _1 ^\infty \frac{e^{-t} t^{\sigma - 1}}{1 - e^{-t}} dt \leq \int _1 ^\infty \frac{e^{-t} t^{c - 1}}{1 - e^{-t}} dt \leq \int _0 ^\infty \frac{e^{-t} t^{c - 1}}{1 - e^{-t}} dt = \Gamma(c)\zeta(c).
\end{equation*}
	 This shows that the integral converges uniformly in every strip $1 + \delta \leq \sigma \leq c$, where $\delta > 0$, and therefore represents an analytic function in every such strip, hence also in the half-plane $\sigma > 1$. Therefore, by analytic continuation, our identity holds for all $s$ with $\sigma > 1$.
\end{proof}


\begin{theorem}
	Let $C = C _1 + C_2 + C_3$ denote a positively oriented contour as shown in Figure~\ref{fig:ContourC}
\begin{figure}[!htb]
\begin{minipage}[c]{0.5\textwidth}
\centering
\begin{equation*}
\begin{aligned}
	C_1 \colon \bcbr{c, \infty} &\to \field{C}, \; C_1(r) = r e^{i \pi} \\
	C_2 \colon \cbr{-\pi, \pi} &\to \field{C}, \; C_2(\theta) = c e^{i \theta} \\
	C_3 \colon \bcbr{c, \infty} &\to \field{C}, \; C_3(r) = r e^{-i \pi} \\
\end{aligned}
\end{equation*}
\end{minipage}
\begin{minipage}[c]{0.5\textwidth}
\raggedleft
\begin{tikzpicture}
	[decoration={
		markings,
		mark=at position 0.75cm with {\arrow[line width=1pt]{>}},
		mark=at position 3cm with {\arrow[line width=1pt]{>}},
		mark=at position 5cm with {\arrow[line width=1pt]{>}},
		mark=at position 7cm with {\arrow[line width=1pt]{>}}
	}]
	% The axes
	\draw[help lines, very thin, ->] (-3,0) -- (2,0) coordinate (xaxis);
	\draw[help lines, very thin, ->] (0,-2) -- (0,2) coordinate (yaxis);
	\node[below, gray] at (xaxis) {$\Re(z)$};
	\node[left, gray] at (yaxis) {$\Im(z)$};

	% The path
	\path[draw,line width=0.6pt,postaction=decorate] (2.3:-3) -- +(2.3,0) arc (-170:170:0.7) -- +(-2.3,0);
	% The labels

	\node[below right] {$0$};

	\draw[xshift=0cm] (0,0) node[circle,fill,inner sep=1.2pt](a){};
	\draw[xshift=0.7cm] (0,0) node[circle,fill,inner sep=1.2pt](a){};
	\node at (1.4,0.3) {$c < 2 \pi$};
	\node at (-1.7,-0.5) {$C_1$};
	\node at (0.8,-0.7) {$C_2$};
	\node at (-1.7,0.5) {$C_3$};
\end{tikzpicture}
\end{minipage}
\caption{Contour C}
\label{fig:ContourC}
\end{figure}
	looping around the negative real axis and enclosing the point z = 0. The function
\begin{equation}
\begin{aligned}
	&I\colon \field{C} \to \field{C},\\
	&I(s) = \frac{1}{2 \pi i} \int _C \frac{z^{s - 1} e^z}{1 - e^z} dz
\end{aligned}
\end{equation}
	is entire and
\begin{equation}
	\zeta(s) = \Gamma(1 - s) I(s) \textit{ for } \sigma > 1.
\end{equation}
\end{theorem}
\begin{proof}
	As the integrand
\begin{equation*}
\begin{aligned}
	&f_z \colon \field{C} \to \field{C},\\
	&f_z(s) = \frac{z^{s - 1} e^z}{1 - e^z}
\end{aligned}
\end{equation*}
	is an entire function of $s$ for $z \in \{ z \in \field{C} \colon \abs{\Im(z)} < 2 \pi \land z \neq 0 \}$ we only have to show that the integrals along $C_1$ and $C_3$ converge uniformly on every compact disk $D_R = \{ z \in \field{C} \colon \abs{z} \leq R \}$ for $R \in \field{R}_{> 0}$. Along $C_1$ for $r \geq 1$ we have
\begin{equation*}
	\abs{z^{s - 1}} = r^{\sigma - 1} \abs{e^{-\pi i (\sigma - 1 + i t)}} = r^{\sigma - 1} e^{\pi t} \leq r^{R - 1} e^{\pi R}
\end{equation*}
	and analogously along $C_3$ for $r \geq 1$ we get
\begin{equation*}
	\abs{z^{s - 1}} = r^{\sigma - 1} \abs{e^{\pi i (\sigma - 1 + i t)}} = r^{\sigma - 1} e^{-\pi t} \leq r^{R - 1} e^{\pi R}.
\end{equation*}
	Hence since $e^r - 1 > \frac{e^r}{2}$ for $r \geq \log(2)$ either on $C_1$ or $C_3$ we have
\begin{equation*}
	\abs{\frac{z^{s - 1} e^z}{1 - e^z} \frac{dz}{dr}} \leq \frac{r^{R - 1} e^{\pi R} e^{-r}}{1 - e^{-r}} = \frac{r^{R - 1} e^{\pi R}}{e^{r} - 1} \leq A(R) r^{R - 1} e^{-r}
\end{equation*}
	for $r \geq 1$ and a constant $A(R) \in \field{R}_{> 0}$ only depending on $R$. As
\begin{equation*}
	\forall c, R \in \field{R}_{> 0} \colon \int _c ^\infty r^{R - 1} e^{-r} dr < \infty
\end{equation*}
	we have shown that $C_1$ and $C_3$ converge uniformly on every disk $D_R$. \\
	For the second claim along $C_1$ and $C_3$ we have 
\begin{equation*}
	g(z) \coloneqq \frac{e^z}{1 - e^z} = \frac{e^{r e^{\pm \pi i}}}{1 - e^{r e^{\pm \pi i}}} = \frac{e^{-r}}{1 - e^{-r}} = g(-r).
\end{equation*}
	We write
\begin{equation*}
\begin{aligned}
	2 \pi i I_c(s)
	&= \left(\int _{C_1} + \int _{C_2} + \int _{C_3}\right) z^{s - 1} g(z) dz \\
	&\begin{aligned}
		= \int _\infty ^c r^{s - 1} e^{-\pi i s} g(-r) dr
		&+ i \int _{-\pi} ^{\pi} c^{s - 1} e^{(s - 1) i \theta} c e^{i \theta} g(c e^\theta) d\theta \\ 
		&+ \int _{c} ^\infty r^{s - 1} e^{-\pi i s} g(-r) dr
	\end{aligned} \\
	&= 2 i \sin(\pi s) \int _c ^\infty r^{s - 1} g(-r) dr + i c^s \int _c ^\infty e^{i s \theta} g(c e^{i \theta}) d\theta.
\end{aligned}
\end{equation*}
	Once we have divided $2i$ on both sides we get
\begin{equation*}
	\pi I_c(s) = \sin(\pi s)I_{1, c}(s) + I_{2, c}(s).
\end{equation*}
	Now we let $c \to 0$. For $I_{1, c}$ previous theorem gives us
\begin{equation*}
	\lim\limits _{c \to 0} I_{1, c}(s) = \int _0 ^\infty \frac{r^{s - 1} e^r}{1 - e^r} dr = \Gamma(s)\zeta(s)
\end{equation*}
	For $I_{2, c}$ we have to keep in mind that $g(z)$ is holomorphic for $\abs{z} < 2 \pi$ except a simple pole at $z = 0$. Thus $z g(z)$ is holomorphic anywhere inside $\abs{z} = c < 2 \pi$ and is bounded there for $z = c < 2 \pi$. So we get $\abs{g(z)} < \frac{B}{\abs{z}}$ for some constant $B \in \field{R}_{> 0}$ and obtain the approximation
\begin{equation*}
	\abs{I_{2, c}(s)} \leq \frac{c^\sigma}{2} \int _{-\pi} ^{\pi} e^{-t \sigma} \frac{B}{c} d\theta \leq B e^{\pi \abs{t}} c^{\sigma - 1}.
\end{equation*}
	If $\sigma > 1$ then $\lim\limits _{c \to 0} I_{2, c}(s) = 0$. Hence we get
\begin{equation*}
	\pi I(s) = \sin(\pi s) \Gamma(s)\zeta(s).
\end{equation*}
	By multiplying both sides with $\Gamma(1 - s)$ and plugging in Euler's reflection formula we finally obtain
\begin{equation*}
	\zeta(s) = \Gamma(1 - s) I(s).
\end{equation*}
\end{proof}


\begin{definition}
	As $\Gamma(1 - s)$ is analytic for $\sigma \geq 1, s \neq 1$ and $I(s)$ is entire we can meromorphically continue $\zeta(s)$ and define 
\begin{equation*}
	\zeta(s) = \Gamma(1 - s) I(s) \textit{ for } \sigma \leq 1, s \neq 1.
\end{equation*}
\end{definition}


\begin{theorem}
	We have $\zeta \in \mathcal{H}(\field{C} \mysetminus \fbr{1})$ with a simple pole in $s = 1$ and its residue $Res\cbr{\zeta(s), 1} = 1$.
\end{theorem}
\begin{proof}
	The only singularities that come into question are the poles of $\Gamma(1 - s)$ as $I(s)$ is entire, but we already know that $\zeta(s)$ is analytic for $\sigma > 1$ which leaves $s = 1$ the single possible pole. \\
	To prove the above statement, assume $s = 1$. The integrals along $C_1$ and $C_3$ take the same values and cancel each other out leaving
\begin{equation*}
\begin{aligned}
	I(1) 
	&= \frac{1}{2 \pi i} \int _{C_2} \frac{z^{0} e^z}{1 - e^z} dz = \frac{1}{2 \pi i} Res\cbr{\frac{e^z}{1 - e^z}, 1} \\
	&= \lim\limits _{z \to 0} z \frac{e^z}{1 - e^z} = \lim\limits _{z \to 0} \frac{z}{e^{-z} - 1} = -\lim\limits _{z \to 0} \frac{1}{e^{-z}} = -1.
\end{aligned}
\end{equation*}
	Now we can compute the residue
\begin{equation*}
\begin{aligned}	
	Res\cbr{\zeta(s), 1} 
	&= \lim\limits _{s \to 1} \cbr{s - 1} \zeta(s) \\
	&= -\lim\limits _{s \to 1} \cbr{1 - s} \Gamma(1 - s) I(s) \\
	&= -I(1) \lim\limits _{s \to 1} \Gamma(2 - s) = \Gamma(1) = 1.
\end{aligned}
\end{equation*}
\end{proof}


\section{The Functional Equation}
The functional equation gives a new powerful closed expression of $\zeta$ and allows us to make further assertions about the occurrence of zeros in the whole complex plane.


\begin{lemma}
	Let $r \in \cbbr{0, \pi}$ and $U = \field{C} \mysetminus \fbr{2 n \pi i \colon n \in \field{Z}}$. If
\begin{equation*}
	S(r) = \fbr{z \in \field{C} \colon \forall n \in \field{Z} \colon \abs{z - 2 n \pi i} \geq r} \subset U
\end{equation*}
	and
\begin{equation*}
	g \colon U \to \field{C}, \; g(z) = \frac{e^z}{1 - e^z}
\end{equation*}
	then $g(z)$ is bounded such that
\begin{equation*}
	\exists A(r) \in \field{R}^+ \; \forall z \in S(r) \colon \abs{g(z)} \leq A(r).
\end{equation*}
\end{lemma}
\begin{proof}
	Let $z = x + i y \in \field{C}$. First we assume $\abs{x} \leq 1$. Consider the punctured rectangle
\begin{equation*}
	Q(r) = \fbr{z \in \field{C} \colon \abs{x} \leq 1, \abs{y} \leq \pi, \abs{z} \geq r} \subset S(r)
\end{equation*}
	As $g$ is continuous and $Q(r)$ is compact in $U$ we know that $g \big\vert _{Q(r)}$ has to be bounded by a $B(r) \in \field{R}^+$. The same holds for $g$ in the infinite strip $S(r) \cap \fbr{z \in \field{C} \colon \abs{x} \leq 1}$ because of its $2 \pi i$-periodicity. \\
	Now suppose $\abs{x} > 1$. We have
\begin{equation*}
	\abs{g(z)} = \abs{\frac{e^z}{1 - e^z}} = \frac{e^x}{\abs{1 - e^z}} \leq \frac{e^x}{\abs{1 - e^x}}.
\end{equation*}
	For $x > 1$ we get
\begin{equation*}
	\abs{g(z)} = \frac{e^x}{e^x - 1} = \frac{1}{1 - e^{-x}} \leq \frac{1}{1 - e^{-1}} = \frac{e}{e - 1}
\end{equation*}
	and analogously for $x < -1$ 
\begin{equation*}
	\abs{g(z)} = \frac{e^x}{1 - e^x} \leq \frac{1}{1 - e^{-x}} \leq \frac{1}{1 - e^{-1}} = \frac{e}{e - 1}
\end{equation*}
	and therefore
\begin{equation*}
	\forall z \in S(r) \colon \abs{g(z)} \leq \max \fbr{\frac{e}{e - 1}, B(r)} \eqqcolon A(r).
\end{equation*}
\end{proof}


\begin{theorem}
	We find that
\begin{equation}
	\zeta(1 - s) = \frac{2}{\cbr{2 \pi}^s} \Gamma(s) \cos(\frac{\pi s}{2}) \zeta(s) \textit{ for } \sigma > 1
\end{equation}
	or equivalently
\begin{equation}
	\zeta(s) = 2\cbr{2 \pi}^s-1 \Gamma(1 - s) \sin(\frac{\pi s}{2}) \zeta(1 - s) \textit{ for } \sigma < 0.
\end{equation}
	This is the so called functional equation of $\zeta$. By meromorphic continuation the domain is extendable to $\field{C} \mysetminus \fbr{1}$.
\end{theorem}
\begin{proof}
	Let $C_N = C_{1,N} + C_{2,N} + C_{3,N} + C_{4,N}$ denote a negatively orientated contour for $N \in \field{N}$ as shown in Figure~\ref{fig:ContourCN},
\begin{figure}[!htb]
\begin{minipage}[c]{0.55\textwidth}
\centering
\begin{equation*}
\begin{aligned}
	C_{1,N} \colon \bbr{c, R(N)} &\to \field{C}, \; C_{1,N}(r) = r e^{i \pi} \\
	C_{2,N} \colon \cbr{-\pi, \pi} &\to \field{C}, \; C_{2,N}(\theta) = c e^{i \theta} \\
	C_{3,N} \colon \bbr{c, R(N)} &\to \field{C}, \; C_{3,N}(r) = r e^{-i \pi} \\
	C_{4,N} \colon \cbr{-\pi, \pi} &\to \field{C}, \; C_{4,N}(\theta) = R(N) e^{i \theta}
\end{aligned}
\end{equation*}
\end{minipage}
\begin{minipage}[c]{0.45\textwidth}
\raggedleft
\begin{tikzpicture}
	[decoration={
		markings,
		mark=at position 0.5cm with {\arrow[line width=1pt]{>}},
		mark=at position 2cm with {\arrow[line width=1pt]{>}},
		mark=at position 4cm with {\arrow[line width=1pt]{>}},
		mark=at position 6cm with {\arrow[line width=1pt]{>}},
		mark=at position 8.5cm with {\arrow[line width=1pt]{>}},
		mark=at position 11.5cm with {\arrow[line width=1pt]{>}},
		mark=at position 14.5cm with {\arrow[line width=1pt]{>}},
		mark=at position 17.5cm with {\arrow[line width=1pt]{>}}
	}]
	% The axes
	\draw[help lines, very thin, ->] (-2.5,0) -- (2.5,0) coordinate (xaxis);
	\draw[help lines, very thin, ->] (0,-2.5) -- (0,2.5) coordinate (yaxis);
	\node[below, gray] at (xaxis) {$\Re(z)$};
	\node[left, gray] at (yaxis) {$\Im(z)$};
	% The path
	\path[draw,line width=0.6pt,postaction=decorate] (3.6:-2) -- +(1.3,0) arc (-170:170:0.7) -- +(-1.3,0) arc (176.5:-176.5:2);
	% The labels
	\node[below right] {$0$};
	\draw[xshift=0cm] (0,0) node[circle,fill,inner sep=1.2pt](a){};
	\draw[xshift=0.7cm] (0,0) node[circle,fill,inner sep=1.2pt](a){};
	\draw[xshift=2cm] (0,0) node[circle,fill,inner sep=1.2pt](a){};
	\node at (0.95,0.3) {$c$};
	\node at (2.65,0.3) {$R(N)$};
	\node at (-1.3,-0.5) {$C_{1,N}$};
	\node at (0.9,-0.8) {$C_{2,N}$};
	\node at (-1.3,0.5) {$C_{3,N}$};
	\node at (1.7,1.7) {$C_{4,N}$};
\end{tikzpicture}
\end{minipage}
\caption{Contour $C_N$}
\label{fig:ContourCN}
\end{figure}
	where $0 < c < \pi$ describes the inner and $R(N) = 2(N + 1) \pi$ the outer radius of our two circles $C_{2,N}$ and $C_{4,N}$. Consider
\begin{equation*}
	I_N(s) = \frac{1}{2 \pi i} \int _{C_N} \frac{z^{s - 1} e^z}{1 - e^z} dz.
\end{equation*}
	We first want to prove that $\lim\limits_{N \to \infty} I_N(s) = I(s)$ for $\sigma < 0$. As
\begin{equation*}
\begin{aligned}
	\lim\limits_{N \to \infty} I_N(s)
		&= \lim\limits_{N \to \infty} \frac{1}{2 \pi i} \cbr{\int_{C_{1,N}} + \int_{C_{2,N}} + \int_{C_{3,N}} + \int_{C_{4,N}}} \frac{z^{s - 1} e^z}{1 - e^z} dz \\
		&= \frac{1}{2 \pi i} \cbr{\int_{C_{1}} + \int_{C_{2}} + \int_{C_{3}} + \lim\limits_{N \to \infty}\int_{C_{4,N}}} \frac{z^{s - 1} e^z}{1 - e^z} dz
\end{aligned}
\end{equation*}
	it suffices to show that the last integral vanishes. We know that
\begin{equation*}
	\abs{z^{s - 1}} = \abs{z^{\sigma - 1}} = \abs{R(N)^{\sigma - 1} e^{-t \sigma}} \leq R(N)^{\sigma -1} e^{\pi \abs{t}}
\end{equation*}
	and from the previous lemma we deduce
\begin{equation*}
	\abs{\frac{e^z}{1 - e^z}} \leq A(\pi)
\end{equation*}
	which leads to
\begin{equation*}
\begin{aligned}	
	\abs{\int_{C_{4,N}} \frac{z^{s - 1} e^z}{1 - e^z} dz}
	&\leq \int_{-\pi}^{\pi} \abs{\frac{z^{s - 1} e^z}{1 - e^z} z^\prime} d\theta \\
	&\leq \int_{-\pi}^{\pi} R(N)^\sigma e^{\pi \abs{t}} A(\pi) d\theta \\
	&= 2 \pi R(N)^\sigma e^{\pi \abs{t}} A(\pi)
\end{aligned}
\end{equation*}
	which converges to $0$ for $N \to \infty$ if $\sigma < 0$. Accordingly we can write
\begin{equation*}
	\lim\limits_{N \to \infty} I_N(1 - s) = I(1 - s) \textit{ for } \sigma > 1.
\end{equation*}
	Now we want to determine $I_N(1 - s)$ explicitly by using Cauchy's residue theorem. Let
\begin{equation*}
	R(n) = Res\cbr{\frac{z^{-s} e^z}{1 - e^z}, 2 n \pi i}
\end{equation*}
	where
\begin{equation*}
\begin{aligned}
	R(n) &= \lim\limits_{z \to 2 n \pi i} \cbr{z - 2 n \pi i} \frac{z^{-s} e^z}{1 - e^z} \\
	&= \cbr{2 n \pi i}^{-s} e^{2 n \pi i} \lim\limits_{z \to 2 n \pi i} \frac{z - 2 n \pi i}{1 - e^z} \\
	&= - \cbr{2 n \pi i}^{-s}.
\end{aligned} 
\end{equation*}
	We get
\begin{equation*}
\begin{aligned}	
	I_N(1 - s)
	&= - \sum _{\substack{n = -N \\ n \neq 0}} ^N R(n) = - \sum _{n = 1} ^N R(n) + R(-n) \\
	&= \sum _{n = 1} ^N \cbr{2 n \pi i}^{-s} + \sum _{n = 1} ^N \cbr{-2 n \pi i}^{-s} \\
	&= \cbr{\frac{e^{-\frac{\pi i s}{2}}}{\cbr{2 \pi}^s} + \frac{e^{\frac{\pi i s}{2}}}{\cbr{2 \pi}^s}} \sum _{n = 1} ^N n^{-s} \\
	&= \frac{2}{\cbr{2 \pi}^s} \cos(\frac{\pi s}{2}) \sum _{n = 1} ^N n^{-s}.
\end{aligned}
\end{equation*}
	Finally we obtain
\begin{equation*}
\begin{aligned}
	\zeta(1 - s) &= \Gamma(s) I(1 - s) = \Gamma(s) \lim\limits_{N \to \infty} I_N(1 - s) \\
	&= \frac{2}{\cbr{2 \pi}^s} \Gamma(s) \cos(\frac{\pi s}{2}) \zeta(s)
\end{aligned}
\end{equation*}
	and by substituting s with 1 - s we have
\begin{equation*}
	\zeta(s) = 2\cbr{2 \pi}^{s - 1} \Gamma(1 - s) \sin(\frac{\pi s}{2}) \zeta(1 - s).
\end{equation*}
\end{proof}


\subsection{Trivial Zeros and Critical Strip}


\begin{corollary}
	$\zeta(s)$ has zeros at $s = -2n$ for $n \in \field{N}_0$. These are called the trivial zeros.
\end{corollary}
\begin{proof}
	The functional equation of $\zeta(s)$ becomes zero whenever $\sin(\frac{\pi s}{2})$ vanishes, which in fact happens when $s = -2n$.
\end{proof}


\begin{corollary}
	$\zeta(s)$ can only have nontrivial zeros within the strip $0 \leq \sigma \leq 1$. This region is also known as the critical strip.
\end{corollary}
\begin{proof}
	Consider the second functional equation from above
\begin{equation*}
	\zeta(s) = 2\cbr{2 \pi}^{s-1} \Gamma(1 - s) \sin(\frac{\pi s}{2}) \zeta(1 - s) \textit{ for } \sigma < 0. 
\end{equation*}
	We already know that $\zeta(s)$ has no zeros for $\sigma > 1$. We also have shown that $\Gamma(s)$ doesn't vanish on its domain and the only remaining factor is $\sin(\frac{\pi s}{2})$, which in fact generates the trivial zeros. Thus $\zeta(s)$ has no nontrivial zeros for $\sigma < 0$.
\end{proof}

\appendix
\bibliographystyle{Style}
\bibliography{References}
\nocite{*}

\end{document}